\chapter{Connectedness and Compactness}

\section{Connected Spaces}

\begin{definition}
    Let $X$ be a topological space.
    \begin{enumerate}[label = (\arabic*),itemsep=1pt,topsep=3pt]
        \item A \textui{separation of $X$} is a pair $U,V$ of disjoint and nonempty open subsets of $X$ whose union is $X$. If a separation of $X$ exists, we say $X$ is \textui{disconnected}.
        \item The space $X$ is \textui{connected} if there does not exist a separation of $X$. 
    \end{enumerate}
\end{definition}

\begin{proposition}
    A nonempty topological space $X$ is connected if and only if the only subsets of $X$ that are both open and closed are $\emptyset$ and $X$.
\end{proposition}
    \begin{proof}
        Suppose $X$ is connected. Note that $\emptyset,X$ are by definition open and closed. If there were another subset $A$ such that $A^c = A$, then $X = A \cup A^c$ would be a separation. This is a contradiction, hence $\emptyset,X$ are the only open and closed subsets of $X$.

        For the converse, we consider the contrapositive of the above statement. If $X$ is a disconnected topological space, then there exists nonempty open subsets $U,V$ of $X$ such that $U \cap V = \emptyset$ and $U \cup V = X$. Then $U$ and $V$ are closed and open, as $U^c = V$ and $V^c = U$.
    \end{proof}

\begin{example}
    Consider $[0,1] \cup [2,3] \subseteq \bfR$ with the subspace topology. Then $[0,1] \cup [2,3]$ is not connected since $U = [0,1]$ and $V = [2,3]$ form a separation.
\end{example}

\begin{definition}
    A topological space is \textui{totally disconnected} if the only connected subsets are singletons.
\end{definition}

\begin{example}
    The rationals are totally disconnected. Let $A \subseteq \bfQ$ be a subspace with at least two points. Let $p,q \in A$ and suppose without loss of generality $p < q$. Pick an irrational number $c$ between $p$ and $q$. Define $U:= (\infty,c) \cap A$ and $V:=(c,\infty) \cap A$. These are open, nonempty, and disjoint subsets of $A$. But $U \cup V = A \setminus \{c\} = A$. Since $A$ admits a separation, $A$ is disconnected. As this is true for any subspace of $\bfQ$ with at least two points, we can conclude that $\bfQ$ is totally disconnected.
\end{example}

\begin{example}
    Consider the topological space $(X,\cP(X))$, If $X$ has at least two points, then $X$ is not connected, as $\{p\}$, $X \setminus \{p\}$ form a separation.
\end{example}

\begin{lemma}\label{lemma:subspace-connected}
    Let $Y$ be a subspace of a topological space $X$.
        \begin{enumerate}[label = (\arabic*),itemsep=1pt,topsep=3pt]
            \item $A,B$ is a separation of $Y$;
            \item $Y = A \cup B$, where $A,B \subseteq Y$ are disjoint, nonempty, and $\overline{A} \cap B = A \cap \overline{B} = \emptyset$.
        \end{enumerate}
\end{lemma}
    \begin{proof}
        Let $A,B$ be a separation of $Y$. Then $Y = A \cup B$, where $A,B \subseteq Y$ are disjoint, nonempty, and open in $Y$. We only need to show that $\overline{A} \cap B = A \cap \overline{B} = \emptyset$. Suppose towards contradiction we can find some $p \in \overline{A} \cap B$. Since $B$ is open in $Y$, we can write $B = Y \cap U$, where $U \subseteq X$ is open. Hence $p \in \overline{A} \cap Y \cap U$. Now since $p \in \overline{A}$ and $p \in U$, the set $A \cap U$ is nonempty. Hence we can find some $a \in A \cap U$. But since $A \subseteq Y$, we have $a \in Y \cap U = B$. Thus $a \in A \cap B$, which is a contradiction. By a similar argument, $A \cap \overline{B} = \emptyset$.

        Conversely, we must show that $A$ and $B$ are open in $Y$. Since $\overline{A} \cap B = \emptyset$, note that:
            \begin{equation*}
            \begin{split}
                \overline{A} \cap Y 
                & = \overline{A} \cap (A \cup B) \\
                & = (\overline{A} \cap A ) \cup (\overline{A} \cap B) \\
                & = A \cup \emptyset \\
                & = A.
            \end{split}
            \end{equation*}
        A similar argument will show that $\overline{B} \cap Y = B$. Now since $\overline{A}$ and $\overline{B}$ are closed, by definition this means $A$ and $B$ are closed in $Y$. Since $A = Y \setminus B$ and $B = Y \setminus A$, they are open in $Y$ as well.
    \end{proof}

\begin{example}
    Consider the following subset of $\bfR^2$:
        \begin{equation*}
        \begin{split}
            Y:=\{(x,0) \mid x \in \bfR\} \cup \{(x, {\ts \frac{1}{x}}) \mid x > 0\}
        \end{split}
        \end{equation*}
    Note that $\{(x,0) \mid x \in \bfR\}$ and $\{(x,\frac{1}{x})\mid x > 0\}$ are disjoint and closed in $\bfR^2$. By Lemma~\ref{lemma:subspace-connected}, $\{(x,0) \mid x \in \bfR\}$ and $\{(x, {\ts \frac{1}{x}}) \mid x > 0\}$ form a separation of $Y$.
\end{example}

\begin{lemma}\label{lemma:connected-subspace-inclusion}
    If $C$ and $D$ form a separation of $X$, and if $Y$ is a connected subspace of $X$, then either $Y \subseteq C$ or $Y \subseteq D$. 
\end{lemma}
    \begin{proof}
        Suppose towards contradiction $Y \cap C \neq \emptyset$ and $Y \cap D \neq \emptyset$. Since $C$ and $D$ are open in $X$, then $Y \cap C$ and $Y \cap D$ are open in $Y$. Moreover, their union is $Y$ and their intersection is empty. Therefore, $Y \cap C$ and $Y \cap D$ form a separation of $Y$, which is a contradiction.
    \end{proof}

\begin{theorem}\label{thm:connected-if-shared-points}
    The union of a collection of connected subspaces of $X$ that have a point in common is connected.
\end{theorem}
    \begin{proof}
        Let $\{A_i\}_{i \in I}$ be a collection of connected subspaces of $X$ and let $p \in \cap_{i \in I}A_i$. Suppose towards contradiction that $U,V$ is a separation of $\bigcup_{i \in I}A_i$. Then either $p \in U$ or $p \in V$. Without loss of generality, assume $p \in U$. Since each $A_i$ is connected, it must lie entirely in either $U$ or $V$, and it cannot lie in $V$ because it contains the point $p \in U$. Hence $A_i \subset U$ for all $i$, which means $\bigcup_{i \in I}A_i \subset U$. But this means that $V = \emptyset$, which is a contradiction.
    \end{proof}

\begin{theorem}
    Let $A$ be a connected subspace of $X$. If $A \subseteq B \subseteq \overline{A}$, then $B$ is connected.
\end{theorem}
    \begin{proof}
        Suppose towards contradiction $B = C \cup D$ is a separation of $B$. Then $A \subseteq C$ or $A \subseteq D$; suppose without loss of generality that $A \subseteq C$. Then $\overline{A} \subseteq \overline{C}$. By Lemma~\ref{lemma:subspace-connected}, $\overline{C} \cap D = \emptyset$. Therefore $B \cap D = \emptyset$, as $B \subseteq \overline{C}$. This contradicts $D$ being a nonempty subset of $B$.
    \end{proof}

\begin{example}
    The \textit{Topologist's sine curve} is the set:
        \begin{equation*}
        \begin{split}
            S := \{(x,\sin {\ts \frac{1}{x}}) \mid 0 < x \leq 1\} \cup \{(0,y) \mid y \in [-1,1]\}.
        \end{split}
        \end{equation*}
    Note that the closure of $\{(x,\sin {\ts \frac{1}{x}}) \mid 0 < x \leq 1\}$ is $S$. Since $\{(x,\sin {\ts \frac{1}{x}}) \mid 0 < x \leq 1\}$ is connected, $S$ is connected.
\end{example}

\begin{theorem}
The image of a connected space under a continuous function is connected.
\end{theorem}
    \begin{proof}
        Let $X$ be connected and $f:X \rightarrow Y$ continuous. If $U$ and $V$ form a separation of $f(X)$, then $f^{-1}(U)$ and $f^{-1}(V)$ form a separation of $X$, which is a contradiction.
    \end{proof}

\begin{theorem}
    A finite Cartesian product of connected spaces is connected
\end{theorem}
    \begin{proof}
        We will only prove the theorem for the product of two connected spaces $X$ and $Y$, as the finite case follows by induction. Choose a point $(a,b) \in X \times Y$. The set $X \times \{b\}$ is connected since it is homeomorphic to $X$. For each $x_0 \in X$, the set $\{x_0\} \times Y$ is connected since it is homeomorphic to $Y$. Let $T_{x_0} := (X \times \{b\}) \cup (\{x_0\} \times Y)$. Since $(x_0,b) \in (X \times \{b\}) \cap (\{x_0\} \times Y)$, by Theorem~\ref{thm:connected-if-shared-points} $T_{x_0}$ is connected. Now $\bigcup_{x_0 \in X}T_{x_0}$ is connected, as $(a,b) \in \bigcap_{x_0 \in X}T_{x_0}$. But $\bigcup_{x_0 \in X}T_{x_0} = X \times Y$, hence $X \times Y$ is connected.
    \end{proof}

\section{Connected Subspaces of $\bfR$}

\begin{theorem}
    A set $E \subseteq \bfR$ is connected if and only if whenever $a < c < b$ with $a,b$, then $c \in E$.
\end{theorem}
    {\color{red} \begin{proof}

    \end{proof}}

\begin{theorem}[Intermediate Value Theorem]
    Let $X$ be a connected topological space and $f:X \rightarrow \bfR$ a continuous map. If $a,b \in X$ and $l$ is a point in $\bfR$ lying between $f(a)$ and $f(b)$, then there exists $c \in X$ such that $f(c) = l$.
\end{theorem}
    \begin{proof}
        Suppose such a $c$ does not exist. Let $A := (-\infty,l) \cap f(X)$ and $B:= (l,\infty) \cap f(X)$. Clearly $A$ and $B$ are disjoint, and moreover $A \cup B = f(X) \setminus \{l\} = f(X)$. This means $A$ and $B$ for a separation of $f(X)$. But $f(X)$ is connected as $f$ is continuous \textemdash a contradiction.
    \end{proof}

\section{Path Connectedness}

\begin{definition}
    For $x,y \in X$, a \textui{path} in $X$ from $x$ to $y$ is a continuous function $\gamma:[0,1] \rightarrow X$ such that $\gamma(0) = x$ and $\gamma(1) = y$. The space $X$ is called \textui{path connected} if every pair of points can be joined by a path.
\end{definition}

\begin{example}
    Any convex subset of a topological vector space is path-connected via a straight line.
\end{example}

\begin{example}
    Let $K$ be a finite subset of $\bfR^n$. Then $\bfR^n \setminus K$ is path connected. This is not true for $n=1$.
\end{example}

\begin{proposition}
    If $X$ is path connected and $f:X \rightarrow Y$ is continuous, then $f(X)$ is path connected.
\end{proposition}
    \begin{proof}
        Let $y_1,y_2 \in f(X)$. Then there exists $x_1,x_2 \in X$ such that $y_1 = f(x_1)$ and $y_2 = f(x_2)$. Since $X$ is path connected, there exists a continuous map $\gamma:[0,1] \rightarrow X$ such that $\gamma(0) = x_1$ and $\gamma(1) = x_2$. Since $f$ and $\gamma$ are continuous, the composition $f \circ \gamma$ is continuous. Hence $f \circ \gamma:[0,1] \rightarrow Y$ is a continuous map such that $(f \circ \gamma)(0) = f(x_1) = y_1$ and $(f \circ \gamma)(1) = f(x_2) = y_2$. Thus $f(X)$ is path connected.
    \end{proof}

\section{Compactness}
\begin{definition}
    Let $X$ be a topological space. A \textui{cover} of $X$ is a collection $\{U_i\}_{i \in I} \subset X$ such that $\bigcup_{i \in I}U_i = X$. If $U_i$ is open for all $i \in I$, we say $\{U_i\}_{i \in I}$ is an \textui{open cover} of $X$. A \textui{subcover} of $\{U_i\}_{i \in I} \subset X$ is a sub-collection of $\{U_i\}_{i \in I}$ whose union is $X$. A \textui{finite subcover} is a finite subcollection of $\{U_i\}_{i \in I}$ that still covers $X$.
\end{definition}

\begin{definition}
    A topological space $X$ is \textui{compact} if every open cover of $X$ admits a finite subcover.
\end{definition}

\begin{example}
    The set $\{(n-1,n+1) \mid n \in \bfN\}$ is an open cover of $\bfR$ which does not admit a finite subcover. Hence $\bfR$ is \textit{not} compact.
\end{example}

\begin{example}
    The set $\left\{ \left( \frac{1}{n+1}, \frac{1}{n} \mid n \in \bfN \right) \right\}$ is an open cover of $(0,1)$ which does not admit a finite subcover.
\end{example}

\begin{example}
    Finite topological spaces are compact. If $X = \{x_1,...,x_n\}$ and $\{U_i\}_{i \in I}$ is an open cover of $X$, then for each $x_i \in X$ pick $U_{x_i} \in \{U_i\}_{i \in I}$ such that $x \in U_{x_i}$. Then $\bigcup_{i=1}^n U_{x_i}$ is a finite subcover.
\end{example}

\begin{example}
    The topological space $(X,\cP(X))$ is not compact. Indeed, $\bigcup_{x \in X}\{x\}$ is an open cover which does not admit a finite subcover.
\end{example}

\begin{proposition}
    Let $Y$ be a subspace of a topological space $X$. Then $Y$ is compact if and only if every collection $\{U_i\}_{i \in I}$ of open sets in $X$ that contains $Y$ has a finite subcollection $\{U_1,...,U_n\}$ also containing $Y$.
\end{proposition}
    {\color{red} \begin{proof}
        
    \end{proof}}

\begin{proposition}
    Every closed subspace of a compact space is compact.
\end{proposition}
    \begin{proof}
        Let $X$ be a compact topological space and let $Y \subseteq X$ be closed. Let $\{U_i\}_{i \in I}$ be an open cover of $Y$. Then $\{U_i\}_{i\in I} \cup Y^c$ is an open cover of $X$. Since $X$ is compact, we can find a finite subcover $\{U_1,...,U_n\} \cup Y^c$. This gives $Y \subseteq \bigcup_{i = 1}^n U_i$, which means $Y$ is compact.
    \end{proof}

\begin{lemma}[A point and a compact set can be separated]
    Let $X$ be a Hausdorff topological space, $K \subseteq X$ a compact subset, and $x_0 \not\in K$. Then there exists open sets $U$ and $V$ such that $x_0 \in U$, $K \subseteq V$, and $U \cap V = \emptyset$.
\end{lemma}
    \begin{proof}
        Since $X$ is Hausdorff, for every $k \in K$ there exists open sets $U_k$ and $V_k$ such that $x_0 \in U_k$, $y \in V_k$, and $U_k \cap V_k = \emptyset$. Then $K \subseteq \bigcup_{k \in K}V_k$. Since $K$ is compact, there exists a finite subcover; i.e., $K \subseteq \bigcup_{i = 1}^n V_{k_i}$. Since each $U_{k_i}$ is a neighborhood of $x_0$ which is disjoint from $V_{k_i}$, it must be the case that $\left( \bigcap_{i = 1}^n U_{k_i} \right) \cap \left( \bigcup_{i = 1}^n V_{k_i} \right) = \emptyset$. This establishes the claim.
    \end{proof}

\begin{proposition}
    A compact subset of a Hausdorff space is closed.
\end{proposition}
    \begin{proof}
        Let $X$ be Hausdorff and $K \subseteq X$ compact. Let $x_0 \in K^c$. By Lemma~\ref{} there exists open sets $U$ and $V$ such that $x_0 \in U$, $K \subseteq V$, and $U \cap V = \emptyset$. But $U \cap K = \emptyset$ implies $U \subseteq K^c$. Thus $x_0 \in U \subseteq K^c$, and since $x_0$ is arbitrary this establishes that $K$ is closed.
    \end{proof}

\begin{example}
    Consider the three point set $\{a,b,c\}$ with the topology:
        \begin{equation*}
        \begin{split}
            T:=\{\emptyset,\{b\},\{a,b\},\{b,c\},\{a,b,c\}\}.
        \end{split}
        \end{equation*}
    Then $\{b\}$ is compact but not closed. 
\end{example}

\begin{proposition}
    The continuous image of a compact set is compact.
\end{proposition}
    \begin{proof}
        Let $f:X \rightarrow Y$ be continuous. Let $\{U_i\}_{i \in I}$ be a collection of open sets in $Y$ that cover $f(X)$. Then $\bigcup_{i \in I}f^{-1}(U_i)$ is an open cover of $X$. Since $X$ is compact, there exists a finite subcover $\{f^{-1}(U_1),...,f^{-1}(U_n)\}$. Hence $\{U_1,...,U_n\}$ is a finite subcover of $f(X)$.
    \end{proof}

\begin{proposition}
    Let $f:X \rightarrow Y$ be a bijective continuous map. If $X$ is compact and $Y$ is Hausdorff, then $f$ is a homeomorphism.
\end{proposition}
    \begin{proof}
        If $C \subseteq X$ is closed, then $C$ is compact. Since $f(C) \subseteq Y$ is compact, $f(C)$ is closed. Thus $f$ is a closed map, which implies it is a homeomorphism.
    \end{proof}

\begin{corollary}
    If $X$ is compact, $Y$ is Hausdorff, and $p:X \rightarrow Y$ is a continuous surjection, then $p$ is a quotient map.
\end{corollary}
